% PTIT-style preliminary pages. Do not reuse identities or statements from sample group.
\chapter*{LỜI CAM ĐOAN}
\addcontentsline{toc}{chapter}{Lời cam đoan}
Nhóm sinh viên thực hiện đề tài gồm Dương Ngọc Phú (N23DCDK055), Phạm Thế Vỹ (N23DCCN205), Vũ Văn Kiên (N23DCDT029) và Lê Hoàng Nguyên Quân (N23DCDT042), thuộc lớp E23CQCE01-N, dưới sự hướng dẫn của TS. Lê Ngọc Hiếu.

Chúng em xác nhận rằng các nội dung mô tả quy trình, dữ liệu và kết quả thực nghiệm trong báo cáo được trình bày nhất quán với dữ liệu, chương trình thực nghiệm và bảng kết quả. Các công trình tham khảo được trích dẫn; các giới hạn về mẫu, thiết kế kiểm định và khả năng tái lập được công khai để người đọc có thể đánh giá phạm vi kết luận.

Nhóm sinh viên chịu trách nhiệm tự kiểm tra tính chính xác và tính nguyên gốc của bản báo cáo trước khi nộp theo quy định của Học viện. Việc ký tên xác nhận do chính các thành viên thực hiện khi hoàn thiện hồ sơ nộp.

\vspace{1.8cm}
\begin{flushright}
TP. Hồ Chí Minh, năm 2026\\[10pt]
Nhóm sinh viên thực hiện
\end{flushright}

\clearpage
\chapter*{LỜI CẢM ƠN}
\addcontentsline{toc}{chapter}{Lời cảm ơn}
Nhóm sinh viên trân trọng cảm ơn giảng viên hướng dẫn TS. Lê Ngọc Hiếu đã đồng hành trong quá trình thực hiện đề tài nghiên cứu khoa học về Ensemble Learning cho phát hiện xâm nhập mạng.

Chúng em cảm ơn các thầy cô của Khoa Công nghệ Thông tin 2, Học viện Công nghệ Bưu chính Viễn thông, vì nền tảng kiến thức và môi trường học thuật hỗ trợ quá trình học tập, nghiên cứu. Báo cáo được thực hiện dựa trên những dữ liệu và nguồn tham khảo có ghi nhận, với mong muốn làm rõ chất lượng thực nghiệm cũng như những hạn chế khi ứng dụng các mô hình học máy cho NIDS.

Nhóm mong nhận được những ý kiến nhận xét và phản biện để tiếp tục hoàn thiện nghiên cứu.

\vspace{1.8cm}
\begin{flushright}
TP. Hồ Chí Minh, năm 2026\\[10pt]
Nhóm sinh viên thực hiện
\end{flushright}
